\documentclass[11pt]{article}

\usepackage[review]{acl}

\usepackage{times}
\usepackage{latexsym}
\usepackage[T1]{fontenc}
\usepackage[utf8]{inputenc}
\usepackage{microtype}
\usepackage{inconsolata}
\usepackage{graphicx}
\usepackage{booktabs}

\title{Non-Exchangeability in KV-Cache Relay: From Fixed Pipelines to Recursive, Model-Driven Delegation}

\author{Anonymous ACL submission}

\begin{document}
\maketitle

\begin{abstract}
Relaying key-value (KV) cache between agents in multi-agent LLM systems
rests on an untested assumption: that some dimension of the cache --
which layer, which channel, which question's content -- is interchangeable
for the purpose of compression or reuse. We test this directly with a
causal audit: we substitute the relayed cache with zeroed, random, or
another question's real content and measure whether downstream accuracy
tracks the substitution. We run this audit across three structurally
distinct multi-agent topologies of increasing realism -- a fixed
sequential pipeline, a static fan-in decomposition, and, our central
contribution, a \emph{dynamic, model-driven} delegation loop in the style
of Recursive Language Models (RLMs) \citep{rlm2026}, where the model
itself writes and executes code to decide what to delegate and when. The
same three-tier ordering -- real content $>$ wrong-but-real content $>$ no
real content -- holds, statistically significantly, on every topology we
test, including the self-directed one, which is to our knowledge the
first causal audit of content-specific KV relay under model-driven rather
than predetermined delegation. The fan-in topology gives the cleanest
result in the paper ($p=0.0001$ for real vs.\ wrong-but-real content). On
the recursive topology, a second signal -- how often the model commits to
an answer within budget, not only what it answers -- resolves a
comparison final-answer accuracy alone cannot. We ground this progression
in a fixed-pipeline foundation with two further headline results: relayed
KV carries real, content-specific information rather than merely being
non-empty, and a calibration procedure's own confidence does not
predict -- and actively misdirects on -- which of two model architectures
is safe to compress. We benchmark four published KV-relay efficiency
techniques against this audit, report where each transfers and where it
fails, and report honestly where our evidence is strongest and where it
is not: a diagnosed non-determinism in the recursive prototype, since
fixed and confirmed reproducible, and an accuracy-over-text-relay
comparison that remains statistically unresolved even where the causal
content-dependence claim is not.
\end{abstract}

\begin{figure*}[t]
\centering
\includegraphics[width=0.85\textwidth]{figures/fig4_topology_ladder.png}
\caption{The paper in one figure. We causally audit whether relayed KV
cache carries real, content-specific information by substituting it with
zeroed, random, or mismatched content and measuring downstream accuracy.
The same three-tier ordering -- real $>$ mismatched $>$ zeroed/random --
holds across three increasingly realistic multi-agent topologies, including
one the model constructs for itself turn by turn. See
Section~\ref{sec:results} for the full numbers behind every panel.}
\label{fig:teaser}
\end{figure*}

\section{Introduction}
\label{sec:intro}

Relaying key-value cache between agents in a multi-agent LLM system rests
on an assumption that is rarely stated and, to our knowledge, never
directly tested: that some dimension of the cache -- which transformer
layer, which channel within a head, which question's content -- can be
treated as generic or interchangeable for the purpose of compression,
reconstruction, or reuse. Prior audits of this assumption, ours included
in earlier form, fix the multi-agent structure itself: a predetermined
sequence or decomposition of who talks to whom, decided in advance by the
system's designer, not by the model. Recursive Language Models
\citep{rlm2026} describe a materially different regime already seeing
real use: a root model writes and executes code in a live Python REPL to
decide, for itself, what to delegate to a sub-call and when, rather than
following a fixed script. RLM's own sub-call return channel is
text -- a decoded string the root model may or may not even read. This
paper asks whether the same non-exchangeability result holds when that
return channel is real KV cache instead, and, more fundamentally, whether
\emph{any} causal content-dependence claim about relayed KV survives the
move from a designer-fixed topology to one the model itself constructs on
the fly. No existing audit of KV/latent reuse -- ours or the published
work we compare against in Section~\ref{sec:related} -- has tested this.

We build the claim up through three structurally distinct topologies,
each testing more of the design space than the last (Figure~\ref{fig:teaser}).

\paragraph{Contributions.}
\begin{itemize}
\item Within a fixed sequential pipeline, we test three axes of the
exchangeability assumption -- depth, position, and content-identity --
and find all three false, with content-identity established by direct
causal intervention rather than inferred from accuracy alone.
\item We show this generalizes to a static fan-in decomposition, where it
produces the cleanest three-tier separation in the paper.
\item Central to this paper, we show it generalizes further to a
dynamic, model-driven delegation loop built in RLM's own style, with KV
cache substituted for RLM's decoded-text channel -- including one
comparison that final-answer accuracy cannot resolve but a complementary
behavioral signal (how often the model commits within its turn budget)
resolves cleanly.
\item We diagnose and, where possible, correct four published KV-relay
efficiency techniques against this audit \citep{kvcomm2025,turboquant2026,polarquant2026,kivi2024}:
an offset-correction method and a rotation-based quantization scheme both
fail from false interchangeability assumptions; a per-channel fix
resolves the latter; layer-selection's cost is architecture-dependent in
a way its own calibration signal does not predict.
\item We report honestly where our newest evidence is less mature: a
diagnosed non-determinism in the recursive prototype, since fixed and
confirmed reproducible, and an accuracy-over-text-relay comparison that
remains unresolved even where content-dependence is not.
\end{itemize}

Four recent papers anticipate pieces of this argument from different
angles \citep{pitfallskv2026,rethinkredundancy2026,adak2026,canonicalmerge2026}
(Section~\ref{sec:related}), but none tests causal content-dependence
under model-driven delegation -- the dimension we extend all four along.

\section{Related Work}
\label{sec:related}

\paragraph{Cross-agent KV-cache relay.} KVCOMM \citep{kvcomm2025} corrects
KV-cache ``offset drift'' when reusing cache across prefix contexts;
LatentMAS \citep{latentmas2026} relays autoregressively-generated latent
``thoughts'' via KV-cache working memory. We differ in relaying real,
decoded reasoning text through a fixed sequential chain, and in
evaluating on multi-hop QA \citep{hotpotqa2018}, which neither tested on.

\paragraph{Compressing the relayed cache.} We adopt layer-wise selection
(dropping low-importance layers, reconstructing via \texttt{zeros},
\texttt{nearest}, or \texttt{interpolate}) and per-head quantization.
Orthogonal Backfill \citep{orthobackfill2026} injects a low-rank residual
instead of discarding outright; we considered but did not implement it
(see Limitations). KVTC \citep{kvtc2026} applies transform coding for up
to 20$\times$ compression, a third paradigm distinct from uniform
quantization and rotation.

\paragraph{Rotation-based quantization interacts badly with RoPE-encoded
keys.} A rotation scheme in the style of TurboQuant \citep{turboquant2026}
and PolarQuant \citep{polarquant2026} -- distinct papers sharing only two
authors -- produced severely out-of-distribution output on our real model
(Figure~\ref{fig:b_int4_journey}). We traced this to RoPE: K is cached
\emph{after} RoPE mixes channel pairs by a position-dependent phase, and a
RoPE-agnostic rotation on top disrupts that structure; an isolating
ablation (rotating only V) confirmed the interaction was with K
specifically, not a generic quantization artifact.

\paragraph{RoPE-aware KV cache quantization.} KIVI \citep{kivi2024}
quantizes keys per-channel and values per-token rather than by rotating;
RotateKV \citep{rotatekv2025} rotates \emph{before} RoPE instead; KVQuant
\citep{kvquant2024} similarly identifies post-RoPE key quantization as
the harder case. Our scoped, KIVI-inspired fix (per-channel keys only)
resolves \texttt{B\_int4}'s collapse -- 0.0\%/0.0\% to 52.0\%/64.1\% F1 at
$n=100$, matching our strongest layer-selection config (McNemar
$p=0.86$) at higher compression, no calibration profile required.

\paragraph{Auditing whether KV reuse does what it claims.}
\citet{kvreusefails2026} shows reuse strategies can silently corrupt an
LLM judge's comparisons even when end-task accuracy looks stable.
\citet{latentcommpay2026} proposes the zeroed/random/mismatched
substitution methodology we adopt directly (Section~\ref{sec:method}),
finding content-specificity is \emph{system}-dependent across three
published systems. We extend this along two axes: cross-architecture
generalization, and delegation-structure generalization -- whether the
audit still finds content-dependence once the topology itself is no
longer fixed. Our own extensions (Section~\ref{sec:results}) arrive at
the same qualitative conclusion while varying different axes -- two
independent audits converging on ``it depends'' is more informative than
either alone.

\paragraph{Model-driven delegation and recursive relay.} Recursive
Language Models \citep{rlm2026} let a root model delegate via a live
Python REPL rather than a fixed topology, reading a sub-call's answer
back as decoded text. LatentMAS and CanonicalMerge
\citep{canonicalmerge2026} relay something richer than text, but both
over a system-fixed topology. CanonicalMerge is close to our own fan-in
topology (near-identical models, a closely related task) but reports only
downstream accuracy, with no causal-intervention test. We build a
genuinely dynamic delegation loop in RLM's own style, substitute real KV
cache for its text channel, and apply our causal-audit methodology to
it -- to our knowledge, no prior work tests whether relayed content is
causally necessary once delegation is no longer system-fixed.

\paragraph{Layer redundancy is architecture-dependent.} ``No Free Swap''
\citep{nofreeswap2026} finds layer interchangeability depends on model
and evaluation protocol, consistent with our own cross-architecture gap
(Section~\ref{sec:results}). ``Rethinking Layer Redundancy''
\citep{rethinkredundancy2026} argues redundancy is a joint function of
model and calibration objective; we test the most direct proxy -- whether
the calibration signal's own confidence predicts safety
(Section~\ref{sec:analysis}) -- and it does not. AdaK \citep{adak2026}
generalizes adaptive KV-budget estimation across our same two
architectures by estimating budget per-instance rather than committing to
a fixed calibration-time ranking -- a positive instance of our principle,
not a counterexample. ``The Pitfalls of KV Cache Compression''
\citep{pitfallskv2026} makes a general version of the point we
investigate concretely: our contribution is a causally-grounded,
cross-axis demonstration within one system, validated via direct audit
rather than inferred from compression metrics alone.

\section{Method}
\label{sec:method}

\subsection{Pipeline}
We study a sequential three-agent pipeline -- Reasoner, Verifier,
Finalizer -- answering HotpotQA \citep{hotpotqa2018} (distractor
configuration) questions, where each item bundles its own ten candidate
passages and requires no retrieval step; our GSM8K \citep{gsm8k2021}
generalization check uses the same structure. The only thing that varies
between configurations is \emph{how} one agent's output reaches the next.
In \texttt{text\_agent}, the Reasoner's decoded text is passed to the
Verifier as a literal string, re-processed from scratch; in every
KV-relay configuration (\texttt{A} through \texttt{E}), the Reasoner's
key-value cache is passed directly, and the Verifier continues generation
from it instead. \texttt{single\_agent} is a no-relay control. All
configurations share identical prompts with \texttt{text\_agent}, so any
accuracy or latency delta is attributable to the relay mechanism itself.

\subsection{Models}
We evaluate on Qwen2.5-7B-Instruct \citep{qwen2p5_2024} and, to test
generalization across architecture, Mistral-7B-Instruct-v0.3
\citep{mistral7b2023} (distinct layer count, attention configuration,
tokenizer, comparable scale; both decoder-only, grouped-query attention,
RoPE). Qwen3-8B is evaluated on \texttt{A}/\texttt{D} only ($n=50$) to
extend the architecture-dependence and calibration-confidence checks to a
third point. A fourth candidate, Phi-3.5-mini-instruct, was dropped after
isolating a pre-existing LongRoPE limitation unrelated to our pipeline
(Limitations).

\subsection{Relay Configurations}
\texttt{B\_int8}/\texttt{B\_int4} apply uniform per-head min-max
quantization to every transmitted key/value tensor. \texttt{C} applies
calibration-driven layer selection, reconstructing dropped layers via
\texttt{zeros}, \texttt{nearest}, or \texttt{interpolate}. \texttt{D}
combines layer selection with adaptive per-layer-tier quantization
(8-bit/4-bit). \texttt{E} adds an anchor-based offset-correction
mechanism adapted from KVCOMM (a negative result, Section~\ref{sec:results}).
\texttt{B\_int4} initially collapsed to 0.0\% accuracy -- diagnosed as a
quantization/RoPE interaction (Section~\ref{sec:related}); our fix
(\texttt{B\_int4\_kivi}), inspired by KIVI, quantizes keys per-channel
instead of per-head, leaving values unchanged.

\subsection{Statistical Methodology}
For every headline comparison we report exact (binomial-form) McNemar's
test on paired per-example correctness and a paired bootstrap 95\%
confidence interval (10,000 resamples) for the difference in mean
token-F1, plus Monte Carlo-estimated power for non-significant
comparisons, so ``not significant'' and ``confirmed no effect'' are never
conflated.

\subsection{Causal Audit}
\label{sec:causal_audit_method}
To test whether accuracy reflects the specific transmitted content rather
than merely \emph{some} non-empty cache, we substitute the relayed KV
with (i) an all-zero cache, (ii) a moment-matched random cache, and (iii)
another held-out question's real cache (from a pool never in the scored
set), applied to our three strongest configurations (\texttt{A},
\texttt{B\_int8}, \texttt{D}). We treat moment-matched-random and
mismatched-example as the more informative evidence and zero-ablation as
a supplementary sanity floor, since zero-ablation is a known-imperfect
causal baseline that can produce large behavioral changes unrelated to
whether the zeroed content was meaningful.

\subsection{Fan-In Topology}
To test whether non-exchangeability is an artifact of the sequential
chain specifically, we build a static fan-in decomposition: HotpotQA's
ten passages split across two independent child agents, RoPE-shifted and
concatenated before one aggregator attends over both. The causal audit
substitutes \emph{every} child's KV (an early design bug left one
unaudited as a shortcut, see Limitations), matching the sequential
pipeline's full-relay-corruption strength.

\subsection{Dynamic, Model-Driven Delegation (RLM+KV)}
Our central topology extension builds a real, sandboxed Python REPL loop
in RLM's own style \citep{rlm2026}: the root model writes and executes
code turn by turn, deciding what to inspect directly and when to delegate
to a child call, until it calls \texttt{final\_answer(...)}. We implement
two sub-call return channels: \texttt{text} (RLM's own mechanism, decodes
to text) and \texttt{kv} (RoPE-shifts and splices the sub-call's real KV
directly onto the root's cache, so the root never observes it as text).
The causal audit substitutes a sub-call's real KV identically to the
fixed topologies. One genuine limitation: the root retains direct
sandboxed access to the original passages independent of any sub-call, so
corrupting a sub-call's KV does not remove all of the root's information
the way it does in the fixed topologies (see Results and Limitations).

We implement this directly on \texttt{transformers} \citep{transformers2020}
rather than \texttt{dspy.RLM}: DSPy is a text-only orchestration layer
with no exposure to KV cache, so our \texttt{kv} channel is not
implementable on top of it.

\section{Results}
\label{sec:results}

\subsection{Experimental Setup}
Unless noted otherwise, all numbers are Qwen2.5-7B-Instruct on HotpotQA
(distractor configuration), $n=100$, exact token-match accuracy and token-F1,
with paired McNemar significance and bootstrap F1 confidence intervals
(Section~\ref{sec:method}). Cross-architecture and topology-generalization
results use $n=50$, stated explicitly in each subsection.

\subsection{Main Results}
Table~\ref{tab:main} reports every relay configuration on our primary
model/task pair. \texttt{A} (uncompressed KV relay) matches
\texttt{single\_agent}'s accuracy while avoiding \texttt{text\_agent}'s
full re-prefill; \texttt{B\_int8} compresses 2$\times$ at no measurable
accuracy cost; \texttt{B\_int4\_kivi} is the Pareto-best point, matching
\texttt{D}'s accuracy at the smallest relayed cache of any configuration
tested (Figure~\ref{fig:compression_pareto}); \texttt{E} fails
catastrophically. We unpack each of these in turn below.

\begin{table*}[t]
\centering
\small
\begin{tabular}{lcccc}
\toprule
\textbf{Config} & \textbf{Accuracy} & \textbf{F1} & \textbf{Latency} & \textbf{KV/hop} \\
\midrule
single\_agent & 57.0\% & 68.4\% & 1.6s & --- \\
text\_agent & 54.0\% & 68.3\% & 7.2s & --- \\
A & 56.0\% & 69.1\% & 8.6s & 144.35 MB \\
B\_int8 & 57.0\% & 70.5\% & 8.3s & 72.16 MB (2.00x) \\
B\_int4 (original) & 0.0\% & 0.0\% & 39.4s & 38.38 MB (4.00x) \\
\textbf{B\_int4\_kivi (fixed)} & \textbf{52.0\%} & \textbf{64.1\%} & 9.6s & 36.43 MB (3.96x vs.\ A) \\
C & 43.0\% & 57.8\% & 10.0s & 104.16 MB \\
D & 50.0\% & 61.9\% & 10.2s & 37.82 MB (3.82x vs.\ A) \\
E & 9.0\% & 17.4\% & 19.8s & 35.63 MB \\
E\_int8 & 1.0\% & 6.2\% & 22.2s & 49.05 MB \\
\bottomrule
\end{tabular}
\caption{Established results (Qwen2.5-7B-Instruct, HotpotQA, $n=100$).
Ratios are compressed-vs-A's 144.35~MB, except \texttt{B\_int4} and
\texttt{B\_int8}, whose self-referential ratio equals their vs.-A ratio
exactly since neither drops any layers. Real two-values-per-byte nibble
packing (Section~\ref{sec:method}) makes 4-bit configurations genuinely,
not just nominally, smaller than 8-bit ones.}
\label{tab:main}
\end{table*}

\begin{figure}[t]
\centering
\includegraphics[width=\columnwidth]{figures/fig3_compression_pareto.png}
\caption{Compression vs.\ accuracy across every relay condition in
Table~\ref{tab:main}. \texttt{B\_int4}'s collapse shows smaller is not
automatically better; \texttt{B\_int4\_kivi} is the Pareto-best point.}
\label{fig:compression_pareto}
\end{figure}

\subsection{Technique Scorecard}
Table~\ref{tab:scorecard} summarizes, at a glance, how every published or
adapted technique we tested fares once audited across the architectures
we evaluate it on. We report this separately from Table~\ref{tab:main}
because a single accuracy number on one model can look like success while
hiding architecture-dependence that only shows up once a second or third
model is tested -- exactly the failure mode Section~\ref{sec:results}'s
remaining findings characterize in detail.

\begin{table*}[t]
\centering
\footnotesize
\setlength{\tabcolsep}{4pt}
\begin{tabular}{p{2.9cm}p{2.8cm}p{2.8cm}p{2.0cm}p{2.9cm}}
\toprule
\textbf{Technique} & \textbf{Qwen2.5-7B} & \textbf{Mistral-7B-v0.3} & \textbf{Qwen3-8B} & \textbf{Verdict} \\
\midrule
Uniform 8-bit quant.\ (\texttt{B\_int8}) & Indistinguishable from uncompressed ($p=1.0$) & Indistinguishable from uncompressed (0 discordant) & not tested & \textbf{Free -- use by default} \\
Naive uniform 4-bit quant.\ (\texttt{B\_int4}) & Total collapse (0.0\%/0.0\%) & not tested (root cause is pipeline-independent) & not tested & \textbf{Broken -- do not use} \\
Rotation-based int4 (TurboQuant / PolarQuant) & Out-of-distribution garbage; traced to RoPE-key interaction & not tested & not tested & \textbf{Broken on keys} \\
Per-channel int4 fix (\texttt{B\_int4\_kivi}, KIVI-inspired) & Matches \texttt{D} ($p=0.86$); 3.96x compression & Matches uncompressed on full causal-audit ladder & not tested & \textbf{Recommended 4-bit fix} \\
Calibration-driven layer selection (\texttt{C}/\texttt{D}) & Mild, not-yet-significant cost ($p=0.34$) & Severe, significant collapse ($-30.5$ F1, $p=0.0001$) & No measurable cost ($p=1.0$) & \textbf{Architecture-dependent -- audit per model} \\
KVCOMM-style offset correction (\texttt{E}) & Catastrophic failure ($p\approx0.00003$ vs.\ \texttt{D}) & not tested (dead end already diagnosed) & not tested & \textbf{Fails -- do not use} \\
\bottomrule
\end{tabular}
\caption{Technique scorecard across every architecture we tested each
technique on. ``Not tested'' is reported explicitly rather than implied
by omission -- several cells reflect a deliberate scoping decision
(Limitations), not a null result.}
\label{tab:scorecard}
\end{table*}

\paragraph{Uniform 8-bit quantization is a free win.} \texttt{A} vs.\
\texttt{B\_int8}: McNemar $p=1.0$, only 3 discordant examples out of 100.
Confirmed on Mistral-7B-Instruct-v0.3 ($n=50$): 0 discordant examples --
identical per-example outcomes. This is not an underpowered null result;
the effect is tight enough at this sample size to be a genuine
near-zero-cost finding.

\paragraph{KVCOMM-style offset correction fails catastrophically, and not
because of a decoding confound.} \texttt{D} vs.\ \texttt{E}: McNemar
$p\approx0.00003$, both before and after an unrelated repetition-penalty
fix (Limitations). Manual inspection of \texttt{E}'s raw output shows
duplicated phrases and hallucinated content, not merely a
lower-quality-but-coherent answer.

\paragraph{\texttt{B\_int4\_kivi} matches our strongest layer-selection
config at genuinely higher compression and lower complexity.}
\texttt{B\_int4\_kivi} (52.0\%/64.1\%) vs.\ \texttt{D} (50.0\%/61.9\%):
McNemar $p=0.86$ -- statistically indistinguishable -- while achieving
3.96x compression versus \texttt{D}'s 3.82x (both vs.\ uncompressed
\texttt{A}), and requiring no per-model calibration profile (byte-accounting
correction history in Appendix~\ref{sec:appendix_verification}).
\texttt{B\_int4\_kivi} also trends below \texttt{A}/\texttt{B\_int8} on
accuracy (52\% vs.\ 56--57\%), not yet statistically confirmed
($p=0.36$--$0.50$) -- an open question, not cost-free. As
Figure~\ref{fig:b_int4_journey} shows, fixing the value side as well does
not improve on this further and is still worse on compression, since
per-token quantization's scale/zero-point overhead vastly exceeds K's
per-channel grouping -- evidence the key/RoPE interaction was the entire
mechanism behind \texttt{B\_int4}'s collapse.

\begin{figure}[t]
\centering
\includegraphics[width=\columnwidth]{figures/fig2_b_int4_journey.png}
\caption{The \texttt{B\_int4} diagnosis-to-fix journey: collapse, a broken
rotation-based attempt, an isolating diagnostic, the working fix
(\texttt{B\_int4\_kivi}), and two value-side variants that did not
improve on it.}
\label{fig:b_int4_journey}
\end{figure}

\subsection{Causal Audit: Does Relayed Content Causally Matter?}
\label{sec:causal_audit_results}

\begin{figure}[t]
\centering
\includegraphics[width=\columnwidth]{figures/fig1_causal_audit.png}
\caption{Causal audit: real $>$ mismatched $>$ zeroed $\approx$ random,
replicated across \texttt{A}, \texttt{D}, and \texttt{B\_int8}. Every
pairwise comparison is statistically significant.}
\label{fig:causal_audit}
\end{figure}

\paragraph{Headline result: relayed KV demonstrably carries specific, real
content, not merely a non-empty cache.} We ran the causal audit
(Section~\ref{sec:causal_audit_method}) on \texttt{A} (uncompressed
relay, sequential chain), $n=50$ (Table~\ref{tab:causal_audit_full}):
every pairwise comparison among real/zeroed/random/mismatched is
significant -- to our knowledge the first direct causal validation of
this kind for a sequential multi-agent KV-relay chain, underwriting every
compression result above as a claim about \emph{preserved information}
rather than an assumption. Qualitatively, \texttt{mismatched}'s wrong
answers are coherent and grammatical -- real linguistic structure
attached to the wrong question -- while \texttt{zeroed} produces garbled
English and \texttt{random} is \emph{more} degraded still (code-fragment
tokens, mixed-language noise) despite matching real activation
statistics: a zero vector is a weak key attention can down-weight, while
realistic random noise resembles genuine signal and actively misdirects
it.

Table~\ref{tab:causal_audit_full} consolidates every subsequent extension
of this audit: five configurations on Qwen2.5-7B-Instruct/HotpotQA, five
on Mistral-7B-Instruct-v0.3, two on Qwen3-8B, one on GSM8K, one on the
C2C cross-architecture bridge, and the two topology extensions central to
this paper (fan-in, RLM+KV). Across every Qwen2.5/HotpotQA configuration
(\texttt{D}, \texttt{B\_int8}, \texttt{C}, \texttt{B\_int4\_kivi}), the
same three-tier ordering replicates with real vs.\ zeroed/random
significant at $p<0.0001$ in every case and real vs.\ mismatched
significant throughout ($p=0.0013$--$0.0192$); a single GSM8K
\citep{gsm8k2021} confirmatory run replicates it more decisively still
($p<0.0002$ throughout), showing the claim is not an artifact of one
task.

\begin{table*}[t]
\centering
\small
\begin{tabular}{llccccc}
\toprule
\textbf{Config} & \textbf{Model / Task / Topology} & \textbf{Real} & \textbf{Zeroed} & \textbf{Random} & \textbf{Mismatched} & \textbf{$n$} \\
\midrule
\multicolumn{7}{l}{\emph{Qwen2.5-7B-Instruct, HotpotQA, sequential chain}} \\
A             & --                           & 50.0/64.0 & 0.0/0.0  & 0.0/0.0  & 28.0/37.1 & 50 \\
D             & --                           & 54.0/60.2 & 0.0/0.0  & 0.0/0.2  & 26.0/32.2 & 50 \\
B\_int8       & --                           & 54.0/68.0 & 0.0/0.0  & 0.0/0.0  & 24.0/34.4 & 50 \\
C             & --                           & 46.0/55.7 & 0.0/0.0  & 0.0/0.0  & 24.0/38.9 & 50 \\
B\_int4\_kivi & --                           & 50.0/61.7 & 0.0/0.0  & 0.0/0.0  & 28.0/34.7 & 50 \\
\midrule
\multicolumn{7}{l}{\emph{Generalization: model, task, and topology}} \\
A             & Mistral-7B-Instruct-v0.3     & 56.0/63.3 & 0.0/0.5  & 0.0/0.3  & 16.0/26.4 & 50 \\
C             & Mistral-7B-Instruct-v0.3     & 20.0/30.3 & 0.0/0.5  & 0.0/0.1  & 14.0/26.1 & 50 \\
D             & Mistral-7B-Instruct-v0.3     & 22.0/32.8 & 0.0/0.5  & 0.0/0.1  & 16.0/29.3 & 50 \\
B\_int8       & Mistral-7B-Instruct-v0.3     & 56.0/63.3 & 0.0/0.5  & 0.0/0.1  & 20.0/29.3 & 50 \\
B\_int4\_kivi & Mistral-7B-Instruct-v0.3     & 54.0/63.2 & 0.0/0.5  & 0.0/0.1  & 16.0/25.0 & 50 \\
A             & Qwen3-8B                     & 54.0/67.2 & 46.0/55.8& 0.0/0.0  & 64.0/74.8 & 50 \\
D             & Qwen3-8B                     & 56.0/71.2 & 46.0/55.8& 2.0/5.0  & 52.0/62.4 & 50 \\
A             & GSM8K (Qwen2.5)              & 90.0/---  & 0.0/---  & 2.0/---  & 28.0/---  & 50 \\
F\_c2c        & C2C bridge (Qwen2.5$\to$Qwen3-8B) & 60.0/70.9 & 0.0/0.0 & 0.0/0.0 & 56.0/70.6 & 50 \\
kv (fan-in)   & Fan-in topology (Qwen2.5)    & 40.0/53.0 & 0.0/0.8  & 0.0/0.0  & 6.0/12.0  & 50 \\
kv (RLM+KV)   & Dynamic delegation (Qwen2.5) & 24.0/32.3 & 8.0/11.2 & 6.0/9.2  & 10.0/15.5 & 50 \\
\bottomrule
\end{tabular}
\caption{Every real/zeroed/random/mismatched causal-audit comparison in
this paper: accuracy/F1 (\%) per condition. \emph{Real} vs.\
\emph{zeroed}/\emph{random} is significant ($p<0.0001$) in every row
except Qwen3-8B, where content-dependence itself is architecture-dependent
(see text). \emph{Real} vs.\ \emph{mismatched} is significant in all rows
except Mistral \texttt{C}/\texttt{D} (underpowered on an already-degraded
layer-selection baseline), Qwen3-8B (a distinct, architecture-dependent
finding), and the C2C bridge (channel presence matters, channel identity
does not -- see text). GSM8K reports accuracy only (no partial-credit F1
metric for this task).}
\label{tab:causal_audit_full}
\end{table*}

\paragraph{Mistral: a consistent, explainable split between quantization
and layer-selection.} \texttt{B\_int4\_kivi} and \texttt{B\_int8}
replicate the full five-comparison ladder as cleanly on Mistral as on
Qwen, but \texttt{C} and \texttt{D} show a more nuanced partial pattern:
mismatched clearly beats zeroed/random, but neither is statistically
distinguishable from mismatched itself ($p=0.51$, $p=0.375$) -- unlike on
Qwen. The likely explanation is not that content-identity stops mattering
on Mistral, but that \texttt{C}/\texttt{D} already perform poorly there
before any auditing (20.0\%/22.0\%, well below Qwen's 46.0\%/54.0\%) --
consistent with Mistral's broader layer-selection fragility, not a
failure to replicate.

\paragraph{Qwen3-8B: content-dependence itself is architecture-dependent.}
Random substitution still causes total, highly significant collapse
($p<0.0001$), but real is \textbf{not} statistically distinguishable from
either zeroed or mismatched (all four comparisons $p=0.30$--$0.79$, F1
CIs include zero). Verified as a real effect, not a bug (no pool leakage;
selective overlap across conditions; full detail in
Appendix~\ref{sec:appendix_verification}). This extends the same
meta-point as the calibration-confidence finding below: even this
paper's own central mechanism varies by architecture, with a
stronger/larger model here robust enough to partially recover from losing
or misdirecting real content while remaining defenseless against
structural noise.

\paragraph{The C2C cross-architecture bridge: channel presence matters,
channel identity does not.} We ran the released C2C fuser checkpoint
\citep{c2c2026} bridging Qwen2.5-7B-Instruct to Qwen3-8B and
causally audited it the same way. Real vs.\ zeroed/random is as decisive
as any comparison in this paper ($p<0.0001$ each), but real vs.\
mismatched is \textbf{not} significant ($p=0.625$), nor is real vs.\
no-communication-at-all ($p=0.754$) -- a genuinely different shape of
result: the bridge requires a real, well-formed projection to be present
at all, but not detectably \emph{which} question produced it. Raw text
rules out a scoring artifact (Appendix~\ref{sec:appendix_verification}).

\subsection{Generalization Across Topologies}
\label{sec:topology_generalization}

\paragraph{Fan-in: the cleanest three-tier result in this paper.} We
built the static fan-in decomposition described in
Section~\ref{sec:method}. Substituting \emph{every} child's KV ($n=50$;
see Limitations for a real experimental-design bug this required fixing)
gives the full three-tier ordering (Table~\ref{tab:causal_audit_full}),
more decisively than the sequential chain: real vs.\ zeroed/random
$p<0.0001$ each; real vs.\ mismatched $p=0.0001$, 18 of 19 discordant
pairs favoring real -- the strongest version of this comparison anywhere
in this paper. Unlike the recursive topology below, this static
decomposition has no alternate pathway for a corrupted child's answer to
route around, so zeroed/random collapse to genuinely near-total zero. One
nuance: mismatched vs.\ zeroed/random is not significant by exact
McNemar (only 3 discordant pairs), but the F1 bootstrap CI excludes zero
for both. A garbage-signature scan of all 200 generations gives a clean,
monotonic qualitative match: 0/50 flagged for real, 4/50 mismatched,
23/50 zeroed, 48/50 random.

\paragraph{Dynamic delegation: content-dependence survives a deliberately
unfavorable stress test.} We built the RLM-style dynamic delegation loop
described in Section~\ref{sec:method}, explicitly the least favorable of
our three topologies for demonstrating anything: our distractor-setting
HotpotQA questions do not need RLM-style decomposition at all, so a
real-model REPL loop constructing its own strategy here solves a harder
coordination problem than the task requires. Consistent with that
framing, \texttt{kv} vs.\ \texttt{text} accuracy remains statistically
unresolved at $n=50$ ($p=1.0$) -- we do not claim an accuracy win here.
The causal audit does resolve, against all three substitutions
($p=0.0078$, $p=0.0039$, $p=0.0156$) -- the full three-tier signature,
confirmed on a topology the model constructs for itself. One caveat:
zeroed/random do not collapse to near-zero the way every other audit in
this paper does, because here uniquely the root retains direct sandboxed
access to the source passages -- corrupting a sub-call's KV does not
remove all of the root's information, so this audit is real and
significant but structurally narrower than the fixed topologies' own.

A second signal, mined from the same transcripts at no additional cost,
gives a more complete confirmation: the model's own \emph{behavior} --
whether it commits to an answer within budget -- tracks the same
ordering, including the one comparison accuracy could not resolve.
Timeout rate is monotonic across conditions (real 22\%, mismatched 42\%,
zeroed 72\%, random 76\%); substituting ``finished within budget'' for
``correct'' in the same McNemar test, every comparison is significant,
including mismatched vs.\ random ($p<0.0001$), the leg accuracy left
unresolved. An earlier version of this run showed diagnosed
non-determinism (Limitations); after a forced-determinism fix, every
number here reproduces bit-identically against an independent rerun.

\paragraph{The same behavioral mechanism also produces a real latency
cost, not a benefit.} Although \texttt{kv} does not win on accuracy
against \texttt{text} above, it is measurably slower: 11.4s vs.\ 9.6s
mean wall-clock time per example ($n=50$), despite splicing only a
child's generated-answer span -- cheap relative to re-encoding that same
answer as text for the root to re-prefill. The gap decomposes into the
same mechanism as the timeout-rate signal above, not a cost in the splice
operation itself: \texttt{kv} both times out more often (22\% vs.\ 14\%,
and a timed-out session costs 2--3$\times$ a finished one, 19.6--23.0s
vs.\ 7.4--9.1s) and, even restricted to sessions that finish within
budget, needs more root turns (6.79 vs.\ 5.95) and more delegation calls
(1.38 vs.\ 1.14 \texttt{llm\_query} calls) before committing to an
answer. This qualifies our efficiency framing elsewhere in this paper:
KV-relay's prefill-avoidance advantage, real in the sequential pipeline
(Table~\ref{tab:main}), does not transfer to a topology where the
relayed unit is a short delegated answer rather than a large shared
context whose re-encoding was itself expensive -- the behavioral cost of
working with less legible content outweighs the small compute saving
here.

\subsection{Analysis}
\label{sec:analysis}

\paragraph{Layer-selection's accuracy cost is architecture-dependent, not
a fixed property of a given compression ratio.} On Mistral ($n=50$), a
near-identical proportion of dropped layers to Qwen's \texttt{D} costs
30.5 F1 points ($p=0.0001$); on Qwen2.5-7B ($n=100$) the cost is smaller
and not-yet-significant ($p=0.34$), though Qwen's most generous
uncertainty bound is still under half of Mistral's confirmed collapse.
Qwen3-8B ($n=50$) sharpens this further: \texttt{D} numerically edges out
\texttt{A}, but this is noise ($p=1.0$) -- a second model where the cost
is indistinguishable from zero, against Mistral's single confirmed
collapse -- evidence for architecture-dependence, not a fixed
``usually costly/free'' property.

\paragraph{A calibration signal's own confidence does not predict
downstream safety, and fails in the wrong direction.} The separation
between calibration-importance scores for kept versus dropped layers:
Qwen's mean is 0.491, Mistral's \emph{larger} at 0.673 -- Mistral's
calibration looks more decisive, yet Mistral collapses harder when that
ranking is acted upon. This rules out a noisier signal as the
explanation: the calibration objective can be equally or more confident
while being less reliable, undetectable from its output alone. Adding
Qwen3-8B (separation 0.591, near-zero cost) gives a trend consistent with
this (Pearson $r=0.619$), but $n=3$ is barely more informative than
$n=2$.

\paragraph{Architecture-dependent layer-selection failures differ in
kind, not only in rate.} \texttt{D}'s incorrect answers differ in
character between models: Qwen's are short, close near-misses; Mistral's
include a long tail of fabricated tangents and repetition loops, not
explained by a decoding confound (ruled out directly;
Appendix~\ref{sec:appendix_verification}).

\paragraph{Donor-question content occasionally, but rarely, bleeds
through under mismatched substitution.} Reviewing incorrect
\texttt{A\_audit\_mismatched} answers against their donor question, we
find two instances (one unambiguous) of the donor's content leaking into
the wrong answer, out of 36 reviewed -- genuine but rare, not the
dominant explanation for the accuracy loss (full review in
Appendix~\ref{sec:appendix_verification}).

\section{Conclusion}
We set out to test an assumption that nearly every KV-relay efficiency
technique makes implicitly: that some dimension of the cache is
interchangeable. It is not. Across depth, position, and content-identity,
and across three increasingly realistic topologies -- a fixed sequential
chain, a static fan-in decomposition, and a dynamic delegation loop the
model constructs for itself -- relayed KV cache behaves as a specific,
non-fungible artifact rather than a generic resource. This holds across
three model families, with one informative exception (Qwen3-8B's partial
robustness to content substitution) that itself supports the broader
claim: non-exchangeability depends on architecture, not a fixed property
of KV-cache relay. Four published techniques land on different sides of
this line for diagnosable reasons -- uniform quantization respects it,
naive rotation and offset-correction do not, and a RoPE-aware per-channel
fix restores what the naive approach breaks. We release this causal-audit
methodology as the more durable contribution: future KV-relay techniques
should be checked against it before their numbers are trusted.

\section*{Limitations}
\label{sec:limitations}

\paragraph{Causal audit scope and remaining topology questions.} The
causal audit now covers every relay condition in this paper except
\texttt{B\_int4} itself, which is uninformative to audit given it already
collapses unaudited. The bleed-through analysis is narrower: a manual
review of 20 of 36 incorrect \texttt{A\_audit\_mismatched} examples by one
annotator, with no automated classifier or inter-annotator agreement
check -- a qualitative, exploratory observation, not a validated rate.
Beyond the three topologies tested here, a genuinely adversarial fan-in
topology \citep{latentagentslie2026} and LatentMAS's layer-wise KV
concatenation \citep{latentmas2026} remain untested against our specific
causal-audit methodology; both are real, cited, not-yet-attempted
extensions, not gaps we found and characterized without naming.

\paragraph{Single-process, single-GPU evaluation.} Every experiment runs
within one Python process on one RTX 4090: no \texttt{KVMessage} is ever
serialized or transmitted over a network. Compression-ratio and
transmitted-byte accounting reflects what a real distributed deployment
would need to send and is meaningful on that basis, but our latency
numbers exclude network transfer time, and smaller relayed caches are not
observed to be faster in this setup, since the model still executes every
transformer layer regardless of what its injected cache contains. This
single-GPU constraint also bounded the sample sizes and the number of
model/config combinations we were able to test.

\paragraph{Statistical power and replication.} Several comparisons --
INT8's near-zero cost, offset correction's catastrophic failure,
Mistral-7B's layer-selection collapse -- are significant at conventional
thresholds even at our evaluation scale. Others are not: at $n=100$, we
cannot yet statistically distinguish Qwen2.5-7B's layer-selection
condition from its uncompressed baseline in isolation (power
$\approx$15\%; roughly 3$\times$ the current sample size would be
needed). Our cross-architecture claim rests on the independently
well-powered Mistral-7B result, not the underpowered Qwen-only
comparison. \texttt{B\_int4\_kivi}'s comparison against \texttt{D} rests
on a single $n=100$ run each, not a multi-seed or replicated evaluation.

\paragraph{Scope: three models (one only partially), mostly one task,
greedy decoding.} We test Qwen2.5-7B-Instruct and Mistral-7B-Instruct-v0.3
on the full suite of configurations, and Qwen3-8B on \texttt{A}/\texttt{D}
only ($n=50$). A fourth candidate, Phi-3.5-mini-instruct, was evaluated
and dropped: after fixing two real bugs in how our pipeline hands decode
off to \texttt{model.generate()} under non-default RoPE scaling, the
model still failed to produce coherent output once the pipeline's
cumulative relayed context crossed roughly 4{,}096 tokens -- the model's
own \texttt{original\_max\_position\_embeddings}. We isolated this to a
pre-existing limitation of the model/library combination itself: plain,
unmodified \texttt{model.generate()}, with no relay or pipeline code
involved at all, degrades into the identical failure signature once
forced past the same threshold. Whether findings -- particularly
layer-selection's architecture-dependence -- extend to substantially
larger or smaller models, mixture-of-experts architectures, or models
outside this decoder-only, grouped-query-attention, RoPE lineage remains
largely untested beyond this one additional data point. All of our own
architecture variation is \emph{between} pipeline runs (the same model
family fills every agent role within a given run); whether a single
pipeline mixing architectures across agent roles is even viable is
untested and, to our knowledge, an open problem in this literature more
broadly -- KVCOMM itself flags this as future work, and we inherit that
gap rather than closing it. Nearly all results are on HotpotQA; we ran
one $n=50$ confirmatory check on GSM8K and both headline findings
replicated cleanly, but this is one run on one additional task family,
not a systematic multi-task study. All results use greedy decoding; we
did not evaluate under sampling or self-consistency.

\paragraph{Nibble-packing is unit-tested, not deployment-tested.} We
implement real two-values-per-byte packing for every \texttt{bits==4}
tensor, verified via round-trip pack/unpack exactness and bit-identical
\texttt{compress}/\texttt{decompress} output with and without packing on
synthetic tensors. This is standalone tensor-level testing, not an
end-to-end GPU rerun with the new packing path in the live pipeline; the
reported accuracy numbers throughout this paper come from pre-packing
runs, and only the compression-ratio figures reflect the packing-enabled
compressor.

\paragraph{Partial reproduction of adapted and considered techniques.}
Our rotation-based quantization covers PolarQuant's core rotate-then-
quantize mechanism but omits TurboQuant's QJL bias-correction term; our
KIVI-inspired fix omits the original method's group-wise windowing and
fp16 residual buffer, implementing full-sequence per-channel quantization
only. We considered but did not implement a fourth layer-reconstruction
strategy, Orthogonal Backfill \citep{orthobackfill2026}: its published
formula requires attention weights this pipeline's \texttt{sdpa} decode
path does not expose, and given our own evidence that naive/simplified
reproductions can actively mislead, we chose not to ship an approximated
version under time pressure.

\paragraph{The recursive topology's evidence is less mature than the rest
of the paper, by design, not by oversight.} Getting a working, coherent
baseline required diagnosing and fixing several real failure modes (a
sub-call's own prompt framing leaking into the receiver's attention when
spliced without isolating the generated answer span; a session-wide
repetition-tracking loss specific to a persistent, turn-continuing cache;
a RoPE position-shift bug specific to splicing an answer-only, prompt-
truncated child KV). Even with those fixed, the \texttt{kv}-vs-\texttt{text}
accuracy comparison remains genuinely unresolved at $n=50$ -- we report
the causal-audit result, which does resolve, as this topology's real
contribution, and leave the accuracy comparison open rather than claim
more than the evidence supports.

\paragraph{A boundary we found by testing where RLM's own motivation
actually lives.} We additionally tested \texttt{kv} vs.\ \texttt{text}
at the context lengths RLM's own motivating regime targets
(2,000--16,000 tokens, beyond HotpotQA's native $\sim$2--3K-token
setting). \texttt{text} decisively outperforms \texttt{kv} at the
default delegation granularity (10 of 10 discordant pairs favoring
\texttt{text}, both lengths, $n=25$ each) -- the opposite direction from
the native-scale comparison above. We hypothesized this reflects
\texttt{kv}'s uncompressed splice diluting relevant signal across a
large delegated batch, and tested whether finer-grained delegation
closes the gap; across two independent samples (batch size 3, $n=25$
each, pooled $n=50$) it does not: \texttt{kv}'s own accuracy is
identical (8.0\%) regardless of batch size, and \texttt{text} continues
to win more discordant pairs in both independent samples (McNemar
$p=0.34$ pooled, F1 CI $[-0.227,+0.039]$) -- directionally consistent,
not formally significant, but sufficient to rule out signal dilution
within a single delegated batch as the explanation.

The real mechanism, found by reusing turn-count and timeout data every
run on this branch already saves (no new inference required): at every
long-context configuration we tested, \texttt{kv} fails to finish within
its turn budget far more often than \texttt{text}, and shrinking the
delegation batch size makes this \emph{worse}, not better, because
covering the same haystack in smaller batches requires proportionally
more delegation calls, and \texttt{kv} was already too slow to resolve
within budget at the easier, larger batch size. Substituting ``finished
within budget'' for ``correct'' in the same paired McNemar test used
throughout this paper: at batch size 20 ($n=25$), $p=0.0005$ (12 of 12
discordant pairs favor \texttt{text} finishing); at batch size 3 (pooled
$n=50$), $p<0.000001$ (32 of 33 discordant pairs favor \texttt{text}).
\texttt{kv}'s mean delegation-call count grows sharply as batch size
shrinks (3.76 at batch size 20 $\to$ 13.24 and 21.04 at batch size 3
across our two independent samples) while \texttt{text}'s stays flat
($\approx$2.4--2.7) --
the opposite of what the dilution hypothesis predicted, and the reason
finer-grained delegation did not rescue \texttt{kv}: it does not reduce
the total work \texttt{kv} needs to do, it only redistributes that work
into more rounds \texttt{kv} cannot complete in time. This is the same
behavioral mechanism (more turns and delegation calls needed to act on
less legible spliced content than on legible decoded text) that explains
\texttt{kv}'s latency cost at native scale above, now additionally
explaining an accuracy deficit at long context: a timeout rate this
lopsided, not a difference in what either channel does once it finishes,
is the proximate cause of \texttt{text} winning the long-context
comparison. We report this boundary, and its mechanism, rather than
leaving the long-context deficit unexplained: near-parity at native
scale does not extend to the longer-context regime motivating the RLM
citation itself, and the reason is a turn-budget asymmetry our own
transcripts make directly visible. Doubling the turn budget (40
$\to$ 80, $n=15$) does not resolve this: \texttt{kv}'s timeout rate is
unchanged (53\% vs.\ 52--76\%), and per-example turn counts are sharply
bimodal -- sessions either resolve in under 35 turns or exhaust the
entire budget with no examples in between, with the latter group
further splitting between near-zero delegation attempts and delegation
loops that never terminate. The deficit is not turn-budget starvation;
whatever makes spliced content harder to act on than decoded text does
not resolve given more time. The causal-audit content-dependence
signature, by contrast, strengthens at this same long-context scale
(\texttt{kv} vs.\ random substitution, completion rate: $n=30$,
$p<0.0001$, 18 of 18 discordant pairs) -- this paper's central thesis
holds even where the accuracy/efficiency comparison does not.

We checked ARR's guidance before deciding whether a separate Ethical
considerations section was warranted: this work uses two long-established
public benchmarks (HotpotQA, GSM8K), openly-released model checkpoints,
no human subjects, no crowdsourcing, and no new data collection -- we are
not aware of a concern specific to this paper beyond what this section
already covers.

\bibliography{references_rlm_focus}
\bibliographystyle{acl_natbib}

\appendix
\section{Implementation Verification and Correction History}
\label{sec:appendix_verification}

This appendix collects the diagnostic and verification detail behind
claims made in the main text, for readers who want to audit our process
rather than only our results.

\paragraph{Qwen3-8B content-dependence verification (Section~\ref{sec:causal_audit_results}).}
We verified the null content-dependence result on Qwen3-8B is real, not a
bug, three ways: no donor-question pool leakage; the zeroed-condition
byte-identical overlap between \texttt{A} and \texttt{D} traces cleanly
to \texttt{D}'s own \texttt{zeros} reconstruction strategy composing
correctly with the audit's own zeroing (not a bug); and real, zeroed, and
mismatched overlap on only 10--24 of 50 examples pairwise, confirming the
substitution takes effect selectively rather than being silently
no-op'd.

\paragraph{C2C raw-text confirmation (Section~\ref{sec:causal_audit_results}).}
Raw generated text confirms the C2C bridge's null real-vs-mismatched
result is not a scoring artifact: zeroed/random produce the same
degenerate-repetition and out-of-vocabulary failure signatures documented
throughout this paper, now observed on a completely different mechanism
(a learned cross-model projector, not quantization or layer dropping) --
the same qualitative texture, not merely the same numbers.

\paragraph{Nibble-packing correction (Section~\ref{sec:results}).} Our
quantized tensors were initially stored as one full \texttt{torch.uint8}
per element regardless of bit-width -- no nibble-packing existed anywhere
in the compressor, so a 4-bit value took the same byte an 8-bit value
did. We caught this, then rather than only correcting the byte
\emph{accounting}, implemented real two-values-per-byte nibble packing
(\texttt{\_pack\_nibbles}/\texttt{\_unpack\_nibbles}), wired into
\texttt{compress()} (pack after quantizing, \texttt{bits==4} only) and
\texttt{decompress()} (unpack before dequantizing). Eleven new unit tests
cover round-trip pack/unpack exactness (including odd-length/padding
edge cases and full-nibble-range coverage) and, most importantly,
bit-identical \texttt{compress}/\texttt{decompress} output with vs.\
without packing -- proving packing is a pure storage-format change with
zero effect on dequantized values. The corrected, packing-backed numbers
land back at the originally-reported figures almost exactly, because the
original formula's arithmetic (0.5 bytes/element for 4-bit) was always
what real packing would produce -- the bug was that nothing in the
compressor actually packed anything.

\paragraph{Fan-in audit design bug (Section~\ref{sec:topology_generalization}).}
An early run of the fan-in causal audit looked like a near-null result
(3/10 EM for \texttt{kv} vs.\ 3/10 for \texttt{kv\_audit\_zeroed}), but
reading the raw predictions showed the misleading aggregate was a real
experimental-design bug, not a null finding: the default
\texttt{causal\_audit\_child\_idx=-1} only corrupted one of two children,
leaving the other's real content intact -- a weaker test than intended.
We fixed this by adding \texttt{causal\_audit\_child\_idx="all"}
(substitute every child), extracted into a directly-testable
\texttt{\_resolve\_audit\_target\_idxs()} static method with five unit
tests. The rerun with \texttt{"all"} is what Section~\ref{sec:topology_generalization}
reports.

\paragraph{RLM+KV non-determinism fix (Section~\ref{sec:topology_generalization}).}
An earlier version of the recursive-delegation run was not run-to-run
deterministic: diffing against an independent rerun on the identical 50
questions found 30\% of \texttt{text}-channel raw generations differing
under greedy decoding, plausibly GPU floating-point non-associativity
cascading through this prototype's long, self-referential, code-writing
decode loop -- unlike the fixed topologies, which are bit-identical
across reruns. We forced deterministic CUDA/cuBLAS/cuDNN kernels
(\texttt{torch.use\_deterministic\_algorithms}, a fixed
\texttt{CUBLAS\_WORKSPACE\_CONFIG} set before any CUDA context exists) and
reran the full five-condition ladder independently: every reported
number -- exact match, F1, timeout rate, and delegation-call count --
reproduces bit-identically against the pre-fix run (0 of 250 raw answers
differ across all five conditions). The one residual nuance: 33 of the
\texttt{random} condition's 50 examples show different intermediate
turn-by-turn wording before both runs independently exhaust their turn
budget on the identical question with the identical outcome -- the
highest-entropy condition's exploratory path is not perfectly pinned
down, but no scored quantity in the paper is affected.

\paragraph{Phi-3.5-mini-instruct diagnosis (Limitations).}
Two real pipeline bugs were found and fixed first: (a) our decode
function's \texttt{position\_offset==0} fast path handed decode off to
\texttt{model.generate()}, which cannot reconstruct true absolute
position across an externally-primed cache under LongRoPE; (b) an
identical unguarded fast path existed in our intermediate-generation
function for Reasoner/Verifier hops, missed by the first fix. After both
fixes, accuracy was still 0.0\%/2.0\%. We isolated the remaining cause
via two scoped GPU diagnostics: (1) our decode function called directly
with a real, chat-formatted, self-consistent donor cache deliberately
crossing the 4,096-token \texttt{original\_max\_position\_embeddings}
threshold mid-decode produced fully coherent output, proving the fixed
code path itself is correct; (2) plain, unmodified
\texttt{model.generate()} with zero pipeline code involved, forced via
\texttt{min\_new\_tokens} past $\sim$4,096 total tokens on a fresh
prompt, degraded into the identical repetition-garbage signature right
around the threshold. (2) is decisive: this is a pre-existing
Phi-3.5-mini + LongRoPE limitation in the installed \texttt{transformers}
version, present even with no relay, injection, or manual loop involved.

\paragraph{Mistral's layer-selection failure character, full detail
(Section~\ref{sec:analysis}).} \texttt{D}'s incorrect answers differ
sharply in character between models. Qwen's are short (mean 18
characters, 2\% exceed 150) and overwhelmingly close, traceable
near-misses. Mistral's are markedly longer (mean 109 characters, 15\%
exceed 150, maximum 1,454) and include a genuine long tail of severe
failures -- fluent but entirely fabricated tangents, degenerate
repetition loops -- absent from Qwen's failure set. We checked an obvious
alternative explanation -- under-penalizing repetition for Mistral
relative to its own tuned defaults -- against each model's published
generation defaults, and find the opposite: Qwen's own default
\emph{is} our \texttt{repetition\_penalty=1.05}, while Mistral's
specifies none, so our setting applies \emph{more} correction for
Mistral, not less. The long-tail pattern persists anyway, weakening a
decoding confound as the explanation and supporting depth-axis
non-exchangeability differing in \emph{character}, not only magnitude,
across architectures.

\paragraph{Long-context timeout failure signatures (Limitations, \S\ref{sec:limitations}).}
Reading the eight timed-out \texttt{kv} transcripts from the
turn-budget-doubling check directly, two distinct, nameable failure
modes account for all of them: four abandon delegation after 1--3
attempts and fall back to reading individual passages one at a time
(\texttt{print(passages[i])}, sequentially, up through index 78 of
$\sim$90 in one case -- arithmetically unable to finish this way within
any reasonable budget); four instead get stuck re-issuing an
\emph{identical} \texttt{llm\_query} call turn after turn, in one case
despite the system's own repetition-detection nudge correctly firing
(\textit{``that is the exact same code as your previous attempts, and it
will produce the same result again''}) -- the loop is detected and
named, but not escaped. Neither failure mode resembles a session that is
gradually making progress and running short on time.

\paragraph{Donor-question bleed-through review, full detail
(Section~\ref{sec:analysis}).} Unlike ``When Latent Agents Lie''
\citep{latentagentslie2026}, which studies an adversarial agent
deliberately substituting hidden state in a fan-in topology, our
substitution is a non-adversarial intervention in a sequential chain:
when the receiving agent is wrong, does its answer reflect the donor's
content, or ordinary confusion on the real question? Manually reviewing
20 of 36 incorrect \texttt{A\_audit\_mismatched} answers against their
donor question, we find one unambiguous instance -- a Kansas university
fight-song question produced ``Ellie Goulding'' in its wrong answer,
traceable only to a donor question about that singer -- and one weaker,
ambiguous case; the remainder show ordinary failure on the real
question's own topic. A genuine but rare phenomenon, not the dominant
explanation for \texttt{A\_audit\_mismatched}'s accuracy loss -- a manual
review of a subset, not an exhaustive classification.

\end{document}
